\documentclass[12pt,a4paper]{article}
\usepackage[utf8]{inputenc}
\usepackage[russian]{babel}
\usepackage{geometry}
\usepackage{multirow}
\usepackage{graphicx}
\usepackage{amsmath}
\usepackage{listings}
\usepackage{xcolor}

\geometry{
    left=30mm,
    right=15mm,
    top=20mm,
    bottom=20mm
}

% Настройка для отображения кода LaTeX
\lstset{
    language=[LaTeX]TeX,
    basicstyle=\small\ttfamily,
    keywordstyle=\color{blue},
    commentstyle=\color{gray},
    stringstyle=\color{red},
    showstringspaces=false,
    breaklines=true,
    frame=single,
    numbers=left,
    numberstyle=\tiny\color{gray},
    literate={а}{{\cyra}}1 {б}{{\cyrb}}1 {в}{{\cyrv}}1 {г}{{\cyrg}}1 {д}{{\cyrd}}1
             {е}{{\cyre}}1 {ё}{{\cyryo}}1 {ж}{{\cyrzh}}1 {з}{{\cyrz}}1 {и}{{\cyri}}1
             {й}{{\cyrishrt}}1 {к}{{\cyrk}}1 {л}{{\cyrl}}1 {м}{{\cyrm}}1 {н}{{\cyrn}}1
             {о}{{\cyro}}1 {п}{{\cyrp}}1 {р}{{\cyrr}}1 {с}{{\cyrs}}1 {т}{{\cyrt}}1
             {у}{{\cyru}}1 {ф}{{\cyrf}}1 {х}{{\cyrh}}1 {ц}{{\cyrc}}1 {ч}{{\cyrch}}1
             {ш}{{\cyrsh}}1 {щ}{{\cyrshch}}1 {ъ}{{\cyrhrdsn}}1 {ы}{{\cyrery}}1
             {ь}{{\cyrsftsn}}1 {э}{{\cyrerev}}1 {ю}{{\cyryu}}1 {я}{{\cyrya}}1
             {А}{{\CYRA}}1 {Б}{{\CYRB}}1 {В}{{\CYRV}}1 {Г}{{\CYRG}}1 {Д}{{\CYRD}}1
             {Е}{{\CYRE}}1 {Ё}{{\CYRYO}}1 {Ж}{{\CYRZH}}1 {З}{{\CYRZ}}1 {И}{{\CYRI}}1
             {Й}{{\CYRISHRT}}1 {К}{{\CYRK}}1 {Л}{{\CYRL}}1 {М}{{\CYRM}}1 {Н}{{\CYRN}}1
             {О}{{\CYRO}}1 {П}{{\CYRP}}1 {Р}{{\CYRR}}1 {С}{{\CYRS}}1 {Т}{{\CYRT}}1
             {У}{{\CYRU}}1 {Ф}{{\CYRF}}1 {Х}{{\CYRH}}1 {Ц}{{\CYRC}}1 {Ч}{{\CYRCH}}1
             {Ш}{{\CYRSH}}1 {Щ}{{\CYRSHCH}}1 {Ъ}{{\CYRHRDSN}}1 {Ы}{{\CYRERY}}1
             {Ь}{{\CYRSFTSN}}1 {Э}{{\CYREREV}}1 {Ю}{{\CYRYU}}1 {Я}{{\CYRYA}}1
             {\_}{{\textunderscore}}1
}

\begin{document}

% Титульный лист
\begin{titlepage}
    \centering
    \vspace*{2cm}
    
    {\Large Министерство науки и высшего образования Российской Федерации\\}
    \vspace{0.5cm}
    {\Large Федеральное государственное образовательное учреждение\\}
    {\Large высшего образования\\}
    \vspace{0.5cm}
    {\Large\bfseries «Российский Университет Дружбы Народом имени Патриса Лумумбы»\\}
    \vspace{1cm}
    {\large Факультет Физико-математических и Естественных наук\\}
    {\large Кафедра Теории Вероятностей и Кибербезопасности\\}
    
    \vspace{3cm}
    
    {\LARGE\bfseries ОТЧЕТ\\}
    \vspace{0.5cm}
    {\Large по лабораторной работе №5\\}
    \vspace{0.5cm}
    {\Large\bfseries «Создание и форматирование таблиц в LaTeX»\\}
    
    \vfill
    
    \begin{flushright}
        Выполнил: студент группы НПМмд02-24\\
        Викторов Е.И.\\
        \vspace{0.5cm}
    \end{flushright}
    
    \vspace{1cm}
    
    {\large Москва, \the\year}
    
\end{titlepage}

% Содержание
\tableofcontents
\newpage

% Основная часть
\section{Цель работы}

Изучить основные принципы создания и форматирования таблиц в системе компьютерной вёрстки LaTeX, освоить использование окружения \texttt{tabular}, команд выравнивания содержимого ячеек и объединения столбцов с помощью команды \texttt{multicolumn}.

\section{Задачи работы}

\begin{enumerate}
    \item Создать простые таблицы с различными типами выравнивания содержимого
    \item Исследовать поведение LaTeX при несоответствии количества элементов в строке и объявленного количества столбцов
    \item Освоить использование команды \texttt{\textbackslash multicolumn} для объединения столбцов
    \item Создать сложные таблицы с комбинированным форматированием
    \item Проанализировать полученные результаты и сформулировать выводы
\end{enumerate}

\section{Теоретическая часть}

\subsection{Основы создания таблиц в LaTeX}

Для создания таблиц в LaTeX используется окружение \texttt{tabular}. Базовый синтаксис:

\begin{verbatim}
\begin{tabular}{спецификация_столбцов}
    содержимое таблицы
\end{tabular}
\end{verbatim}

\subsection{Спецификация столбцов}

Спецификация столбцов определяет количество столбцов и выравнивание содержимого в них:

\begin{itemize}
    \item \texttt{l} — выравнивание по левому краю (left)
    \item \texttt{c} — выравнивание по центру (center)
    \item \texttt{r} — выравнивание по правому краю (right)
    \item \texttt{|} — вертикальная линия между столбцами
\end{itemize}

Например, \texttt{\{|l|c|r|\}} создаёт таблицу с тремя столбцами и вертикальными линиями.

\subsection{Разделители в таблице}

\begin{itemize}
    \item \texttt{\&} — разделитель ячеек в строке
    \item \texttt{\textbackslash\textbackslash} — конец строки
    \item \texttt{\textbackslash hline} — горизонтальная линия
\end{itemize}

\subsection{Команда multicolumn}

Синтаксис команды \texttt{\textbackslash multicolumn}:

\begin{verbatim}
\multicolumn{количество_столбцов}{выравнивание}{содержимое}
\end{verbatim}

Эта команда позволяет объединить несколько столбцов в один.

\section{Практическая часть}

\subsection{Эксперимент 1: Простая таблица}

Создана базовая таблица 3×3 с различными типами выравнивания:

\begin{center}
\begin{tabular}{|l|c|r|}
\hline
Левый & Центр & Правый \\
\hline
Текст & Текст & Текст \\
\hline
123 & 456 & 789 \\
\hline
\end{tabular}
\end{center}

\textbf{Код:}
\begin{verbatim}
\begin{tabular}{|l|c|r|}
\hline
Левый & Центр & Правый \\
\hline
Текст & Текст & Текст \\
\hline
123 & 456 & 789 \\
\hline
\end{tabular}
\end{verbatim}

\subsection{Эксперимент 2: Выравнивание по левому краю}

Таблица с выравниванием всех столбцов по левому краю:

\begin{center}
\begin{tabular}{|l|l|l|}
\hline
Имя & Фамилия & Возраст \\
\hline
Иван & Иванов & 25 \\
Мария & Петрова & 30 \\
Александр & Сидоров & 28 \\
\hline
\end{tabular}
\end{center}

\textbf{Результат:} Все данные выровнены по левому краю, что удобно для текстовой информации.

\subsection{Эксперимент 3: Выравнивание по центру}

Таблица с центральным выравниванием:

\begin{center}
\begin{tabular}{|c|c|c|}
\hline
Товар & Количество & Цена \\
\hline
Яблоки & 10 & 100 \\
Груши & 5 & 150 \\
Апельсины & 8 & 200 \\
\hline
\end{tabular}
\end{center}

\textbf{Результат:} Содержимое всех ячеек расположено по центру, таблица выглядит симметрично.

\subsection{Эксперимент 4: Выравнивание по правому краю}

Таблица с выравниванием по правому краю (удобно для числовых данных):

\begin{center}
\begin{tabular}{|r|r|r|}
\hline
Год & Доход & Расход \\
\hline
2021 & 1000000 & 800000 \\
2022 & 1200000 & 900000 \\
2023 & 1500000 & 1000000 \\
\hline
\end{tabular}
\end{center}

\textbf{Результат:} Числовые данные выровнены по правому краю, что облегчает их сравнение.

\subsection{Эксперимент 5: Комбинированное выравнивание}

Таблица с различными типами выравнивания в разных столбцах:

\begin{center}
\begin{tabular}{|l|c|r|}
\hline
Название (слева) & Категория (центр) & Цена (справа) \\
\hline
Ноутбук & Электроника & 50000 \\
Книга & Литература & 500 \\
Стул & Мебель & 3000 \\
\hline
\end{tabular}
\end{center}

\textbf{Результат:} Каждый столбец имеет оптимальное выравнивание для своего типа данных.

\subsection{Эксперимент 6: Недостаточное количество элементов}

Исследование поведения LaTeX при указании меньшего количества элементов, чем столбцов:

\begin{center}
\begin{tabular}{|l|c|r|r|}
\hline
Столбец 1 & Столбец 2 & Столбец 3 & Столбец 4 \\
\hline
Полная строка & Данные & 100 & 200 \\
\hline
Только два & элемента & & \\
\hline
Один & & & \\
\hline
Все четыре & столбца & заполнены & здесь \\
\hline
\end{tabular}
\end{center}

\textbf{Результат:} LaTeX успешно компилирует таблицу, оставляя пустые ячейки справа. Ошибок не возникает.

\subsection{Эксперимент 7: Избыточное количество элементов}

При попытке указать больше элементов, чем объявлено столбцов, возникает ошибка компиляции:

\texttt{! Extra alignment tab has been changed to \textbackslash cr.}
Это означает, что количество элементов превышает количество объявленных столбцов.

\subsection{Эксперимент 8: Базовое использование multicolumn}

Объединение столбцов для создания заголовка таблицы:

\begin{center}
\begin{tabular}{|l|c|r|}
\hline
\multicolumn{3}{|c|}{\textbf{Заголовок таблицы}} \\
\hline
Имя & Возраст & Город \\
\hline
Иван & 25 & Москва \\
Мария & 30 & Санкт-Петербург \\
\hline
\end{tabular}
\end{center}

\textbf{Код:}
\begin{verbatim}
\multicolumn{3}{|c|}{\textbf{Заголовок таблицы}}
\end{verbatim}

\textbf{Результат:} Три столбца объединены в один с центральным выравниванием.

\subsection{Эксперимент 9: Объединение двух столбцов}

\begin{center}
\begin{tabular}{|l|c|c|r|}
\hline
Продукт & \multicolumn{2}{c|}{Характеристики} & Цена \\
\hline
Телефон & Память & 128 ГБ & 30000 \\
Планшет & Экран & 10 дюймов & 25000 \\
\hline
\end{tabular}
\end{center}

\textbf{Результат:} Два средних столбца объединены для создания общего заголовка.

\subsection{Эксперимент 10: Сложная таблица}

Создание сложной таблицы с несколькими уровнями заголовков:

\begin{center}
\begin{tabular}{|l|c|c|c|c|}
\hline
\multicolumn{5}{|c|}{\textbf{Отчёт за 2023 год}} \\
\hline
\multicolumn{1}{|c|}{Квартал} & \multicolumn{2}{c|}{Доходы} & \multicolumn{2}{c|}{Расходы} \\
\hline
& План & Факт & План & Факт \\
\hline
Q1 & 100 & 110 & 80 & 85 \\
Q2 & 120 & 115 & 90 & 88 \\
Q3 & 130 & 140 & 95 & 100 \\
Q4 & 150 & 155 & 100 & 105 \\
\hline
\multicolumn{1}{|c|}{Итого:} & \multicolumn{2}{c|}{520} & \multicolumn{2}{c|}{378} \\
\hline
\end{tabular}
\end{center}

\textbf{Результат:} Создана многоуровневая таблица с объединёнными ячейками на разных уровнях.

\section{Результаты исследования}

\subsection{Сравнительная таблица типов выравнивания}

\begin{center}
\begin{tabular}{|c|l|l|}
\hline
\textbf{Тип} & \textbf{Обозначение} & \textbf{Применение} \\
\hline
Левое & \texttt{l} & Текстовые данные, названия \\
\hline
Центральное & \texttt{c} & Заголовки, короткие значения \\
\hline
Правое & \texttt{r} & Числовые данные, суммы \\
\hline
\end{tabular}
\end{center}

\subsection{Анализ ошибок}

\begin{center}
\begin{tabular}{|l|l|l|}
\hline
\textbf{Ситуация} & \textbf{Результат} & \textbf{Ошибка} \\
\hline
Мало элементов & Пустые ячейки & Нет \\
\hline
Много элементов & Не компилируется & Да \\
\hline
Правильное количество & Нормальная таблица & Нет \\
\hline
\end{tabular}
\end{center}

\section{Выводы}

По результатам выполнения лабораторной работы были получены следующие выводы:

\begin{enumerate}
    \item \textbf{Выравнивание содержимого:}
    \begin{itemize}
        \item Тип \texttt{l} (левое) оптимален для текстовых данных
        \item Тип \texttt{c} (центральное) подходит для заголовков и симметричного отображения
        \item Тип \texttt{r} (правое) удобен для числовых данных и финансовых показателей
        \item Комбинирование типов выравнивания повышает читаемость таблиц
    \end{itemize}
    
    \item \textbf{Обработка ошибок количества элементов:}
    \begin{itemize}
        \item При недостатке элементов LaTeX создаёт пустые ячейки без ошибок
        \item При избытке элементов возникает ошибка компиляции
        \item Важно строго соблюдать соответствие между объявленным и фактическим количеством столбцов
    \end{itemize}
    
    \item \textbf{Использование multicolumn:}
    \begin{itemize}
        \item Команда \texttt{\textbackslash multicolumn} позволяет гибко объединять столбцы
        \item Можно создавать многоуровневые заголовки
        \item Для каждого объединения можно задать индивидуальное выравнивание
        \item Синтаксис: \texttt{\textbackslash multicolumn\{n\}\{align\}\{text\}}
    \end{itemize}
    
    \item \textbf{Практические рекомендации:}
    \begin{itemize}
        \item Всегда проверять соответствие количества ячеек и столбцов
        \item Использовать подходящее выравнивание для каждого типа данных
        \item Применять \texttt{multicolumn} для улучшения структуры сложных таблиц
        \item Использовать горизонтальные и вертикальные линии для повышения читаемости
    \end{itemize}
    
    \item \textbf{Освоенные навыки:}
    \begin{itemize}
        \item Создание простых и сложных таблиц в LaTeX
        \item Форматирование содержимого ячеек
        \item Объединение столбцов
        \item Отладка ошибок в таблицах
        \item Создание профессионально оформленных документов
    \end{itemize}
\end{enumerate}

\section{Заключение}

В ходе выполнения лабораторной работы были изучены основные принципы создания и форматирования таблиц в системе LaTeX. Освоены команды выравнивания содержимого, исследовано поведение системы при различных ошибках, изучена команда \texttt{\textbackslash multicolumn} для объединения столбцов.

Полученные знания и навыки позволяют создавать таблицы различной сложности для научных статей, отчётов, дипломных работ и другой технической документации. LaTeX обеспечивает высокое качество типографского оформления таблиц и их автоматическое позиционирование в документе.

Система LaTeX показала себя как мощный инструмент для создания профессионально оформленных документов с таблицами любой сложности.

\section{Список использованных источников}

\begin{enumerate}
    \item Львовский С.М. Набор и вёрстка в системе LaTeX. — М.: МЦНМО, 2003. — 448 с.
    \item Котельников И.А., Чеботаев П.З. LaTeX 2ε по-русски. — Новосибирск: Сибирский хронограф, 2004. — 496 с.
    \item The LaTeX Project. Official Documentation. — URL: https://www.latex-project.org/
    \item Overleaf Documentation. Tables. — URL: https://www.overleaf.com/learn/latex/Tables
    \item Wikibooks. LaTeX/Tables. — URL: https://en.wikibooks.org/wiki/LaTeX/Tables
\end{enumerate}

\end{document}
